% Options for packages loaded elsewhere
\PassOptionsToPackage{unicode}{hyperref}
\PassOptionsToPackage{hyphens}{url}
\PassOptionsToPackage{dvipsnames,svgnames,x11names}{xcolor}
%
\documentclass[
]{article}
\usepackage{amsmath,amssymb}
\usepackage{lmodern}
\usepackage{iftex}
\ifPDFTeX
  \usepackage[T1]{fontenc}
  \usepackage[utf8]{inputenc}
  \usepackage{textcomp} % provide euro and other symbols
\else % if luatex or xetex
  \usepackage{unicode-math}
  \defaultfontfeatures{Scale=MatchLowercase}
  \defaultfontfeatures[\rmfamily]{Ligatures=TeX,Scale=1}
\fi
% Use upquote if available, for straight quotes in verbatim environments
\IfFileExists{upquote.sty}{\usepackage{upquote}}{}
\IfFileExists{microtype.sty}{% use microtype if available
  \usepackage[]{microtype}
  \UseMicrotypeSet[protrusion]{basicmath} % disable protrusion for tt fonts
}{}
\makeatletter
\@ifundefined{KOMAClassName}{% if non-KOMA class
  \IfFileExists{parskip.sty}{%
    \usepackage{parskip}
  }{% else
    \setlength{\parindent}{0pt}
    \setlength{\parskip}{6pt plus 2pt minus 1pt}}
}{% if KOMA class
  \KOMAoptions{parskip=half}}
\makeatother
\usepackage{xcolor}
\usepackage{color}
\usepackage{fancyvrb}
\newcommand{\VerbBar}{|}
\newcommand{\VERB}{\Verb[commandchars=\\\{\}]}
\DefineVerbatimEnvironment{Highlighting}{Verbatim}{commandchars=\\\{\}}
% Add ',fontsize=\small' for more characters per line
\newenvironment{Shaded}{}{}
\newcommand{\AlertTok}[1]{\textcolor[rgb]{1.00,0.00,0.00}{\textbf{#1}}}
\newcommand{\AnnotationTok}[1]{\textcolor[rgb]{0.38,0.63,0.69}{\textbf{\textit{#1}}}}
\newcommand{\AttributeTok}[1]{\textcolor[rgb]{0.49,0.56,0.16}{#1}}
\newcommand{\BaseNTok}[1]{\textcolor[rgb]{0.25,0.63,0.44}{#1}}
\newcommand{\BuiltInTok}[1]{\textcolor[rgb]{0.00,0.50,0.00}{#1}}
\newcommand{\CharTok}[1]{\textcolor[rgb]{0.25,0.44,0.63}{#1}}
\newcommand{\CommentTok}[1]{\textcolor[rgb]{0.38,0.63,0.69}{\textit{#1}}}
\newcommand{\CommentVarTok}[1]{\textcolor[rgb]{0.38,0.63,0.69}{\textbf{\textit{#1}}}}
\newcommand{\ConstantTok}[1]{\textcolor[rgb]{0.53,0.00,0.00}{#1}}
\newcommand{\ControlFlowTok}[1]{\textcolor[rgb]{0.00,0.44,0.13}{\textbf{#1}}}
\newcommand{\DataTypeTok}[1]{\textcolor[rgb]{0.56,0.13,0.00}{#1}}
\newcommand{\DecValTok}[1]{\textcolor[rgb]{0.25,0.63,0.44}{#1}}
\newcommand{\DocumentationTok}[1]{\textcolor[rgb]{0.73,0.13,0.13}{\textit{#1}}}
\newcommand{\ErrorTok}[1]{\textcolor[rgb]{1.00,0.00,0.00}{\textbf{#1}}}
\newcommand{\ExtensionTok}[1]{#1}
\newcommand{\FloatTok}[1]{\textcolor[rgb]{0.25,0.63,0.44}{#1}}
\newcommand{\FunctionTok}[1]{\textcolor[rgb]{0.02,0.16,0.49}{#1}}
\newcommand{\ImportTok}[1]{\textcolor[rgb]{0.00,0.50,0.00}{\textbf{#1}}}
\newcommand{\InformationTok}[1]{\textcolor[rgb]{0.38,0.63,0.69}{\textbf{\textit{#1}}}}
\newcommand{\KeywordTok}[1]{\textcolor[rgb]{0.00,0.44,0.13}{\textbf{#1}}}
\newcommand{\NormalTok}[1]{#1}
\newcommand{\OperatorTok}[1]{\textcolor[rgb]{0.40,0.40,0.40}{#1}}
\newcommand{\OtherTok}[1]{\textcolor[rgb]{0.00,0.44,0.13}{#1}}
\newcommand{\PreprocessorTok}[1]{\textcolor[rgb]{0.74,0.48,0.00}{#1}}
\newcommand{\RegionMarkerTok}[1]{#1}
\newcommand{\SpecialCharTok}[1]{\textcolor[rgb]{0.25,0.44,0.63}{#1}}
\newcommand{\SpecialStringTok}[1]{\textcolor[rgb]{0.73,0.40,0.53}{#1}}
\newcommand{\StringTok}[1]{\textcolor[rgb]{0.25,0.44,0.63}{#1}}
\newcommand{\VariableTok}[1]{\textcolor[rgb]{0.10,0.09,0.49}{#1}}
\newcommand{\VerbatimStringTok}[1]{\textcolor[rgb]{0.25,0.44,0.63}{#1}}
\newcommand{\WarningTok}[1]{\textcolor[rgb]{0.38,0.63,0.69}{\textbf{\textit{#1}}}}
\usepackage{longtable,booktabs,array}
\usepackage{calc} % for calculating minipage widths
% Correct order of tables after \paragraph or \subparagraph
\usepackage{etoolbox}
\makeatletter
\patchcmd\longtable{\par}{\if@noskipsec\mbox{}\fi\par}{}{}
\makeatother
% Allow footnotes in longtable head/foot
\IfFileExists{footnotehyper.sty}{\usepackage{footnotehyper}}{\usepackage{footnote}}
\makesavenoteenv{longtable}
\usepackage{graphicx}
\makeatletter
\def\maxwidth{\ifdim\Gin@nat@width>\linewidth\linewidth\else\Gin@nat@width\fi}
\def\maxheight{\ifdim\Gin@nat@height>\textheight\textheight\else\Gin@nat@height\fi}
\makeatother
% Scale images if necessary, so that they will not overflow the page
% margins by default, and it is still possible to overwrite the defaults
% using explicit options in \includegraphics[width, height, ...]{}
\setkeys{Gin}{width=\maxwidth,height=\maxheight,keepaspectratio}
% Set default figure placement to htbp
\makeatletter
\def\fps@figure{htbp}
\makeatother
\setlength{\emergencystretch}{3em} % prevent overfull lines
\providecommand{\tightlist}{%
  \setlength{\itemsep}{0pt}\setlength{\parskip}{0pt}}
\setcounter{secnumdepth}{-\maxdimen} % remove section numbering
% definitions for citeproc citations
\NewDocumentCommand\citeproctext{}{}
\NewDocumentCommand\citeproc{mm}{%
  \begingroup\def\citeproctext{#2}\cite{#1}\endgroup}
\makeatletter
 % allow citations to break across lines
 \let\@cite@ofmt\@firstofone
 % avoid brackets around text for \cite:
 \def\@biblabel#1{}
 \def\@cite#1#2{{#1\if@tempswa , #2\fi}}
\makeatother
\newlength{\cslhangindent}
\setlength{\cslhangindent}{1.5em}
\newlength{\csllabelwidth}
\setlength{\csllabelwidth}{3em}
\newenvironment{CSLReferences}[2] % #1 hanging-indent, #2 entry-spacing
 {\begin{list}{}{%
  \setlength{\itemindent}{0pt}
  \setlength{\leftmargin}{0pt}
  \setlength{\parsep}{0pt}
  % turn on hanging indent if param 1 is 1
  \ifodd #1
   \setlength{\leftmargin}{\cslhangindent}
   \setlength{\itemindent}{-1\cslhangindent}
  \fi
  % set entry spacing
  \setlength{\itemsep}{#2\baselineskip}}}
 {\end{list}}
\usepackage{calc}
\newcommand{\CSLBlock}[1]{\hfill\break\parbox[t]{\linewidth}{\strut\ignorespaces#1\strut}}
\newcommand{\CSLLeftMargin}[1]{\parbox[t]{\csllabelwidth}{\strut#1\strut}}
\newcommand{\CSLRightInline}[1]{\parbox[t]{\linewidth - \csllabelwidth}{\strut#1\strut}}
\newcommand{\CSLIndent}[1]{\hspace{\cslhangindent}#1}
\ifLuaTeX
\usepackage[bidi=basic]{babel}
\else
\usepackage[bidi=default]{babel}
\fi
\babelprovide[main,import]{american}
% get rid of language-specific shorthands (see #6817):
\let\LanguageShortHands\languageshorthands
\def\languageshorthands#1{}
\ifLuaTeX
  \usepackage{selnolig}  % disable illegal ligatures
\fi
\IfFileExists{bookmark.sty}{\usepackage{bookmark}}{\usepackage{hyperref}}
\IfFileExists{xurl.sty}{\usepackage{xurl}}{} % add URL line breaks if available
\urlstyle{same} % disable monospaced font for URLs
\hypersetup{
  pdftitle={NBI-Slurm: Simplified submission of Slurm jobs with energy
saving mode},
  pdfauthor={Andrea Telatin},
  pdflang={en-US},
  colorlinks=true,
  linkcolor={Maroon},
  filecolor={Maroon},
  citecolor={Blue},
  urlcolor={Blue},
  pdfcreator={LaTeX via pandoc}}

\title{NBI-Slurm: Simplified submission of Slurm jobs with energy saving
mode}

\definecolor{c53baa1}{RGB}{83,186,161}
\definecolor{c202826}{RGB}{32,40,38}
\def \rorglobalscale {0.1}
\newcommand{\rorlogo}{%
\begin{tikzpicture}[y=1cm, x=1cm, yscale=\rorglobalscale,xscale=\rorglobalscale, every node/.append style={scale=\rorglobalscale}, inner sep=0pt, outer sep=0pt]
  \begin{scope}[even odd rule,line join=round,miter limit=2.0,shift={(-0.025, 0.0216)}]
    \path[fill=c53baa1,nonzero rule,line join=round,miter limit=2.0] (1.8164, 3.012) -- (1.4954, 2.5204) -- (1.1742, 3.012) -- (1.8164, 3.012) -- cycle;
    \path[fill=c53baa1,nonzero rule,line join=round,miter limit=2.0] (3.1594, 3.012) -- (2.8385, 2.5204) -- (2.5172, 3.012) -- (3.1594, 3.012) -- cycle;
    \path[fill=c53baa1,nonzero rule,line join=round,miter limit=2.0] (1.1742, 0.0669) -- (1.4954, 0.5588) -- (1.8164, 0.0669) -- (1.1742, 0.0669) -- cycle;
    \path[fill=c53baa1,nonzero rule,line join=round,miter limit=2.0] (2.5172, 0.0669) -- (2.8385, 0.5588) -- (3.1594, 0.0669) -- (2.5172, 0.0669) -- cycle;
    \path[fill=c202826,nonzero rule,line join=round,miter limit=2.0] (3.8505, 1.4364).. controls (3.9643, 1.4576) and (4.0508, 1.5081) .. (4.1098, 1.5878).. controls (4.169, 1.6674) and (4.1984, 1.7642) .. (4.1984, 1.8777).. controls (4.1984, 1.9719) and (4.182, 2.0503) .. (4.1495, 2.1132).. controls (4.1169, 2.1762) and (4.0727, 2.2262) .. (4.0174, 2.2635).. controls (3.9621, 2.3006) and (3.8976, 2.3273) .. (3.824, 2.3432).. controls (3.7505, 2.359) and (3.6727, 2.367) .. (3.5909, 2.367) -- (2.9676, 2.367) -- (2.9676, 1.8688).. controls (2.9625, 1.8833) and (2.9572, 1.8976) .. (2.9514, 1.9119).. controls (2.9083, 2.0164) and (2.848, 2.1056) .. (2.7705, 2.1791).. controls (2.6929, 2.2527) and (2.6014, 2.3093) .. (2.495, 2.3487).. controls (2.3889, 2.3881) and (2.2728, 2.408) .. (2.1468, 2.408).. controls (2.0209, 2.408) and (1.905, 2.3881) .. (1.7986, 2.3487).. controls (1.6925, 2.3093) and (1.6007, 2.2527) .. (1.5232, 2.1791).. controls (1.4539, 2.1132) and (1.3983, 2.0346) .. (1.3565, 1.9436).. controls (1.3504, 2.009) and (1.3351, 2.0656) .. (1.3105, 2.1132).. controls (1.2779, 2.1762) and (1.2338, 2.2262) .. (1.1785, 2.2635).. controls (1.1232, 2.3006) and (1.0586, 2.3273) .. (0.985, 2.3432).. controls (0.9115, 2.359) and (0.8337, 2.367) .. (0.7519, 2.367) -- (0.1289, 2.367) -- (0.1289, 0.7562) -- (0.4837, 0.7562) -- (0.4837, 1.4002) -- (0.6588, 1.4002) -- (0.9956, 0.7562) -- (1.4211, 0.7562) -- (1.0118, 1.4364).. controls (1.1255, 1.4576) and (1.2121, 1.5081) .. (1.2711, 1.5878).. controls (1.2737, 1.5915) and (1.2761, 1.5954) .. (1.2787, 1.5991).. controls (1.2782, 1.5867) and (1.2779, 1.5743) .. (1.2779, 1.5616).. controls (1.2779, 1.4327) and (1.2996, 1.3158) .. (1.3428, 1.2113).. controls (1.3859, 1.1068) and (1.4462, 1.0176) .. (1.5237, 0.944).. controls (1.601, 0.8705) and (1.6928, 0.8139) .. (1.7992, 0.7744).. controls (1.9053, 0.735) and (2.0214, 0.7152) .. (2.1474, 0.7152).. controls (2.2733, 0.7152) and (2.3892, 0.735) .. (2.4956, 0.7744).. controls (2.6016, 0.8139) and (2.6935, 0.8705) .. (2.771, 0.944).. controls (2.8482, 1.0176) and (2.9086, 1.1068) .. (2.952, 1.2113).. controls (2.9578, 1.2253) and (2.9631, 1.2398) .. (2.9681, 1.2544) -- (2.9681, 0.7562) -- (3.3229, 0.7562) -- (3.3229, 1.4002) -- (3.4981, 1.4002) -- (3.8349, 0.7562) -- (4.2603, 0.7562) -- (3.8505, 1.4364) -- cycle(0.9628, 1.7777).. controls (0.9438, 1.7534) and (0.92, 1.7357) .. (0.8911, 1.7243).. controls (0.8623, 1.7129) and (0.83, 1.706) .. (0.7945, 1.7039).. controls (0.7588, 1.7015) and (0.7252, 1.7005) .. (0.6932, 1.7005) -- (0.4839, 1.7005) -- (0.4839, 2.0667) -- (0.716, 2.0667).. controls (0.7477, 2.0667) and (0.7805, 2.0643) .. (0.8139, 2.0598).. controls (0.8472, 2.0553) and (0.8768, 2.0466) .. (0.9025, 2.0336).. controls (0.9282, 2.0206) and (0.9496, 2.0021) .. (0.9663, 1.9778).. controls (0.9829, 1.9534) and (0.9914, 1.9209) .. (0.9914, 1.8799).. controls (0.9914, 1.8362) and (0.9819, 1.8021) .. (0.9628, 1.7777) -- cycle(2.6125, 1.3533).. controls (2.5889, 1.2904) and (2.5553, 1.2359) .. (2.5112, 1.1896).. controls (2.4672, 1.1433) and (2.4146, 1.1073) .. (2.3529, 1.0814).. controls (2.2916, 1.0554) and (2.2228, 1.0427) .. (2.1471, 1.0427).. controls (2.0712, 1.0427) and (2.0026, 1.0557) .. (1.9412, 1.0814).. controls (1.8799, 1.107) and (1.8272, 1.1433) .. (1.783, 1.1896).. controls (1.7391, 1.2359) and (1.7052, 1.2904) .. (1.6817, 1.3533).. controls (1.6581, 1.4163) and (1.6465, 1.4856) .. (1.6465, 1.5616).. controls (1.6465, 1.6359) and (1.6581, 1.705) .. (1.6817, 1.7687).. controls (1.7052, 1.8325) and (1.7388, 1.8873) .. (1.783, 1.9336).. controls (1.8269, 1.9799) and (1.8796, 2.0159) .. (1.9412, 2.0418).. controls (2.0026, 2.0675) and (2.0712, 2.0804) .. (2.1471, 2.0804).. controls (2.223, 2.0804) and (2.2916, 2.0675) .. (2.3529, 2.0418).. controls (2.4143, 2.0161) and (2.467, 1.9799) .. (2.5112, 1.9336).. controls (2.5551, 1.8873) and (2.5889, 1.8322) .. (2.6125, 1.7687).. controls (2.636, 1.705) and (2.6477, 1.6359) .. (2.6477, 1.5616).. controls (2.6477, 1.4856) and (2.636, 1.4163) .. (2.6125, 1.3533) -- cycle(3.8015, 1.7777).. controls (3.7825, 1.7534) and (3.7587, 1.7357) .. (3.7298, 1.7243).. controls (3.701, 1.7129) and (3.6687, 1.706) .. (3.6333, 1.7039).. controls (3.5975, 1.7015) and (3.5639, 1.7005) .. (3.5319, 1.7005) -- (3.3226, 1.7005) -- (3.3226, 2.0667) -- (3.5547, 2.0667).. controls (3.5864, 2.0667) and (3.6192, 2.0643) .. (3.6526, 2.0598).. controls (3.6859, 2.0553) and (3.7155, 2.0466) .. (3.7412, 2.0336).. controls (3.7669, 2.0206) and (3.7883, 2.0021) .. (3.805, 1.9778).. controls (3.8216, 1.9534) and (3.8301, 1.9209) .. (3.8301, 1.8799).. controls (3.8301, 1.8362) and (3.8206, 1.8021) .. (3.8015, 1.7777) -- cycle;
  \end{scope}
\end{tikzpicture}
}

%%%%%%%%%%%%%%%%%%%%%%%%%%%%%%%%%%%%%%%%%%%%%%%%%%%%%%%%%%%%%%%%%%%%%%%%
% Authors and Affiliations

\usepackage[affil-it]{authblk}
\usepackage{orcidlink}
%% \renewcommand\Authsep{, }
\setlength{\affilsep}{1em}
\author[1,2%
  %
  ]{Andrea Telatin%
    \,\orcidlink{0000-0001-7619-281X}\,%
    }

\affil[1]{Quadram Institute Bioscience, Norwich, NR4 7UQ, UK%
    \,\protect\href{https://ror.org/04td3ys19}{\protect\rorlogo}\,%
  }
\affil[2]{Center for Microbial Interactions, Norwich Research Park,
Norwich, NR4 7UG, UK%
  }
%%%%%%%%%%%%%%%%%%%%%%%%%%%%%%%%%%%%%%%%%%%%%%%%%%%%%%%%%%%%%%%%%%%%%%%%
\date{18 March 2026}

\begin{document}
\maketitle

\section{Summary}\label{summary}

NBI-Slurm is a Perl package that provides a simplified, user-friendly
interface for submitting and managing jobs on SLURM
(\citeproc{ref-jette2002slurm}{Yoo et al., 2003}) high-performance
computing (HPC) clusters. It offers both a library of Perl modules for
programmatic job management and a suite of command-line tools designed
to reduce the cognitive overhead of SLURM's native interface. Two
distinctive features of NBI-Slurm are TUI applications to view and
cancel jobs, and an energy-aware scheduling mode --- ``eco mode'' ---
that automatically defers flexible jobs to off-peak periods, helping
research institutions reduce their computational carbon footprint
without requiring users to manually plan submission times.

\section{Statement of Need}\label{statement-of-need}

HPC clusters are indispensable in modern research, particularly in the
life sciences where large-scale sequence analyses, genome assemblies,
and statistical models demand resources beyond a desktop workstation.
SLURM has become the dominant workload manager in this space
(\citeproc{ref-slurm_adoption}{Wang et al., 2020}), yet its interface
presents a steep learning curve. Users must learn a verbose
\texttt{sbatch} scripting syntax, understand resource unit conventions
(memory in megabytes, time in \texttt{D-HH:MM:SS} format), manage job
dependencies manually, and repeat boilerplate directives across every
submission script.

Workflow managers such as Snakemake
(\citeproc{ref-molder2021snakemake}{Mölder et al., 2021}) and Nextflow
(\citeproc{ref-di2017nextflow}{Di Tommaso et al., 2017}) address this at
the pipeline level by abstracting SLURM as an execution backend, but
they require users to rewrite their analysis logic inside a
domain-specific language. Many researchers have existing shell scripts
or one-off analyses that do not warrant a full pipeline refactor.
NBI-Slurm occupies a complementary niche: it wraps SLURM's interface
without imposing a workflow model, making it straightforward to submit
individual commands or small batches while retaining access to all SLURM
features through pass-through options.

The \texttt{lsjobs} utility prints a colour-coded, human-readable table
of queued jobs as a static snapshot, offering a more ergonomic
alternative to the raw output of \texttt{squeue}. Its companion tool
\texttt{viewjobs} provides a fully interactive terminal user interface
(TUI) that allows users to browse the live job queue without leaving the
terminal (\autoref{fig:viewjobs}). Users can scroll through jobs with
arrow or Vim keys, sort columns, inspect per-job details, toggle column
visibility, and adjust column widths interactively. Individual jobs can
be selected with \texttt{Space} and multiple selected jobs can be
cancelled in bulk with a single keypress, removing the need to
copy-paste job IDs into \texttt{scancel}.

\begin{figure}
\centering
\includegraphics{viewjobs.png}
\caption{Interactive TUI of \texttt{viewjobs}, showing job navigation,
multi-column display, and bulk-cancel workflow. The image is AI
generated from a real screenshot (Google NanoBanana)
\label{fig:viewjobs}}
\end{figure}

Energy consumption in research computing is a growing concern
(\citeproc{ref-lannelongue2021green}{Lannelongue et al., 2021}). Most
researchers have no practical mechanism to shift flexible jobs to
periods when grid electricity is cheaper or cleaner. NBI-Slurm addresses
this directly with a configurable scheduling module that calculates the
next available low-energy window and injects a \texttt{-\/-begin}
directive into the submission, requiring no change to the underlying
command.

\section{Availability and
Installation}\label{availability-and-installation}

NBI-Slurm is distributed under the MIT licence and is available from
CPAN as \texttt{NBI::Slurm}. Installation requires Perl 5.16 or later
and can be performed with:

\begin{Shaded}
\begin{Highlighting}[]
\ExtensionTok{cpanm}\NormalTok{ NBI::Slurm}
\end{Highlighting}
\end{Shaded}

The source code is hosted at
\url{https://github.com/quadram-institute-bioscience/NBI-Slurm} under
continuous integration. Development has been active since June 2023, and
the module is published to the MetaCPAN repository at
\url{https://metacpan.org/dist/NBI-Slurm}.

\section{Code Structure and
Dependencies}\label{code-structure-and-dependencies}

The package is organised into two layers.

\textbf{Perl module library (\texttt{lib/NBI/}):} Five classes model the
key abstractions:

\begin{itemize}
\tightlist
\item
  \texttt{NBI::Opts} --- encapsulates SLURM resource directives (queue,
  threads, memory, wall-time, email, job arrays, start time). It accepts
  human-friendly inputs such as \texttt{"8GB"} or \texttt{"2h30m"} and
  converts them to SLURM's expected formats.
\item
  \texttt{NBI::Job} --- represents a job to be submitted, holding a
  command (or list of commands) and an \texttt{NBI::Opts} object. The
  \texttt{script()} method generates a complete \texttt{sbatch} script;
  \texttt{run()} submits it and returns the job identifier.
\item
  \texttt{NBI::Queue} --- queries the live SLURM queue via
  \texttt{squeue} and returns a list of \texttt{NBI::QueuedJob} objects,
  optionally filtered by user, status, name, or queue.
\item
  \texttt{NBI::QueuedJob} --- a lightweight data object representing one
  queued job, used by the queue-management tools.
\item
  \texttt{NBI::EcoScheduler} --- implements the energy-aware scheduling
  logic. Given a job's expected duration and a set of configurable
  windows, it finds the next period satisfying a three-tier preference:
  (1) job completes within an eco window and avoids peak hours; (2) job
  starts in an eco window and avoids peak hours but may overrun; (3) job
  starts in an eco window and partially overlaps peak hours. Default
  windows target weekday nights (00:00--06:00) and weekend off-peak
  periods (00:00--07:00, 11:00--16:00), avoiding evening peaks
  (17:00--20:00), all of which are fully configurable with a settings
  file (by default \texttt{\textasciitilde{}/.nbislurm.config}).
\end{itemize}

\textbf{Command-line tools (\texttt{bin/}):}

\begin{longtable}[]{@{}
  >{\raggedright\arraybackslash}p{(\columnwidth - 2\tabcolsep) * \real{0.4000}}
  >{\raggedright\arraybackslash}p{(\columnwidth - 2\tabcolsep) * \real{0.6000}}@{}}
\toprule\noalign{}
\begin{minipage}[b]{\linewidth}\raggedright
Tool
\end{minipage} & \begin{minipage}[b]{\linewidth}\raggedright
Purpose
\end{minipage} \\
\midrule\noalign{}
\endhead
\bottomrule\noalign{}
\endlastfoot
\texttt{runjob} & Submit a command as a SLURM job with resource flags \\
\texttt{lsjobs} & List, filter, and cancel user jobs with coloured
tabular output \\
\texttt{viewjobs} & Interactive terminal UI for job management \\
\texttt{waitjobs} & Block until jobs matching a pattern complete \\
\texttt{whojobs} & Show cluster utilisation grouped by user \\
\texttt{session} & Launch an interactive SLURM session \\
\end{longtable}

Runtime dependencies are deliberately minimal: \texttt{Capture::Tiny}
(\textgreater=0.40), \texttt{JSON::PP}, \texttt{Text::ASCIITable}
(\textgreater=0.22), \texttt{Term::ANSIColor}, \texttt{Storable}, and
\texttt{POSIX}---all either part of the Perl core or widely available on
CPAN.

\section{Documentation}\label{documentation}

Each module is documented with embedded POD (Plain Old Documentation),
rendered on CPAN at \url{https://metacpan.org/dist/NBI-Slurm}. Each
command-line tool provides a \texttt{-\/-help} flag and a manual page
generated from its POD. A user guide with annotated examples is
maintained in the repository's \texttt{README.md}. The test suite
(\texttt{t/}) covers unit behaviour of every module and integration
behaviour of the command-line tools; author-facing tests (\texttt{xt/})
verify POD completeness and coverage. All tests will be able to check
functions even without Slurm. To check the ability to interact with
Slurm, there are optional tests that can be executed with
\texttt{prove\ -lv\ xt/hpc-*.t}.

\section{Example Applications}\label{example-applications}

\textbf{Submitting a parallel job.} A researcher wishing to run a genome
assembler with 18 cores, 64 GB RAM, and a 12-hour wall-time can write:

\begin{Shaded}
\begin{Highlighting}[]
\ExtensionTok{runjob} \AttributeTok{{-}n} \StringTok{"assembly"} \AttributeTok{{-}c}\NormalTok{ 18 }\AttributeTok{{-}m}\NormalTok{ 64 }\AttributeTok{{-}t}\NormalTok{ 12 }\AttributeTok{{-}w}\NormalTok{ ./logs/ }\DataTypeTok{\textbackslash{}}
  \StringTok{"flye {-}{-}nano{-}raw reads.fastq {-}{-}out{-}dir asm"}
\end{Highlighting}
\end{Shaded}

\textbf{Processing a file list as a job array.} To align 200 FASTQ
files, one job per file:

\begin{Shaded}
\begin{Highlighting}[]
\ExtensionTok{runjob} \AttributeTok{{-}n} \StringTok{"align"} \AttributeTok{{-}c}\NormalTok{ 8 }\AttributeTok{{-}m}\NormalTok{ 16 }\AttributeTok{{-}{-}files}\NormalTok{ samples.txt }\DataTypeTok{\textbackslash{}}
  \StringTok{"bwa mem ref.fa \#FILE\# \textgreater{} \#FILE\#.bam"}
\end{Highlighting}
\end{Shaded}

\textbf{Energy-aware deferral.} A long-running but flexible annotation
job can be scheduled for the next eco window automatically. Note that by
default the \emph{eco mode} is enabled, and can be overridden with
\texttt{-\/-no-eco} or setting the economy\_mode=0 in the configuration
file.

\begin{Shaded}
\begin{Highlighting}[]
\ExtensionTok{runjob} \AttributeTok{{-}{-}eco} \AttributeTok{{-}n} \StringTok{"annotate"} \AttributeTok{{-}t}\NormalTok{ 6 }\StringTok{"prokka genome.fa"}
\end{Highlighting}
\end{Shaded}

NBI-Slurm calculates the next suitable window (e.g., the following
night) and adds \texttt{-\/-begin=2026-03-19T00:00:00} to the submission
without any further user action.

\textbf{Programmatic job chaining.} In a Perl analysis script:

\begin{Shaded}
\begin{Highlighting}[]
\FunctionTok{use} \FunctionTok{NBI::Job}\NormalTok{;}
\FunctionTok{use} \FunctionTok{NBI::Opts}\NormalTok{;}

\KeywordTok{my} \DataTypeTok{$opts}\NormalTok{ = }\FunctionTok{NBI::Opts}\NormalTok{{-}\textgreater{}new(}
\NormalTok{    {-}queue =\textgreater{} }\OtherTok{"}\StringTok{long}\OtherTok{"}\NormalTok{, }
\NormalTok{    {-}threads =\textgreater{} }\DecValTok{16}\NormalTok{, }
\NormalTok{    {-}memory =\textgreater{} }\DecValTok{32}\NormalTok{, }
\NormalTok{    {-}}\FunctionTok{time}\NormalTok{ =\textgreater{} }\OtherTok{"}\StringTok{4h}\OtherTok{"}\NormalTok{);}
\KeywordTok{my} \DataTypeTok{$job}\NormalTok{  = }\FunctionTok{NBI::Job}\NormalTok{{-}\textgreater{}new(}
\NormalTok{    {-}name =\textgreater{} }\OtherTok{"}\StringTok{step1}\OtherTok{"}\NormalTok{, }
\NormalTok{    {-}command =\textgreater{} }\OtherTok{"}\StringTok{python align.py}\OtherTok{"}\NormalTok{, }
\NormalTok{    {-}opts =\textgreater{} }\DataTypeTok{$opts}\NormalTok{);}
\KeywordTok{my} \DataTypeTok{$id}\NormalTok{   = }\DataTypeTok{$job}\NormalTok{{-}\textgreater{}}\DataTypeTok{run}\NormalTok{();}

\KeywordTok{my} \DataTypeTok{$opts2}\NormalTok{ = }\FunctionTok{NBI::Opts}\NormalTok{{-}\textgreater{}new(}
\NormalTok{    {-}queue =\textgreater{} }\OtherTok{"}\StringTok{short}\OtherTok{"}\NormalTok{, }
\NormalTok{    {-}threads =\textgreater{} }\DecValTok{4}\NormalTok{, }
\NormalTok{    {-}memory =\textgreater{} }\DecValTok{8}\NormalTok{, }
\NormalTok{    {-}}\FunctionTok{time}\NormalTok{ =\textgreater{} }\OtherTok{"}\StringTok{1h}\OtherTok{"}\NormalTok{);}
\KeywordTok{my} \DataTypeTok{$job2}\NormalTok{  = }\FunctionTok{NBI::Job}\NormalTok{{-}\textgreater{}new(}
\NormalTok{    {-}name =\textgreater{} }\OtherTok{"}\StringTok{step2}\OtherTok{"}\NormalTok{, }
\NormalTok{    {-}command =\textgreater{} }\OtherTok{"}\StringTok{python report.py {-}{-}input results/}\OtherTok{"}\NormalTok{, }
\NormalTok{    {-}opts =\textgreater{} }\DataTypeTok{$opts2}\NormalTok{);}
\DataTypeTok{$job2}\NormalTok{{-}\textgreater{}}\DataTypeTok{opts}\NormalTok{{-}\textgreater{}}\DataTypeTok{dependencies}\NormalTok{([}\DataTypeTok{$id}\NormalTok{]);}
\DataTypeTok{$job2}\NormalTok{{-}\textgreater{}}\DataTypeTok{run}\NormalTok{();}
\end{Highlighting}
\end{Shaded}

\section{Acknowledgements}\label{acknowledgements}

The author gratefully acknowledges the support of the Biotechnology and
Biological Sciences Research Council (BBSRC); this research was funded
by the BBSRC Institute Strategic Programme Food Microbiome and Health
BB/X011054/1 and its constituent project(s) BBS/E/QU/230001B; the BBSRC
Institute Strategic Programme Microbes and Food Safety BB/X011011/1 and
its constituent project(s) BBS/E/QU/230002C; the BBSRC Core Capability
Grant BB/CCG2260/1. This research was also supported by the
infrastructure provided by the CLIMB-BIG-DATA grant MR/T030062/1. The
author thanks colleagues at the Quadram Institute Bioscience for
feedback and field-testing during development, and the GreenDISC working
group and NBI Research Computing for support and discussions.

\section{AI Usage Disclosure}\label{ai-usage-disclosure}

Claude Code (Anthropic) was used during development of NBI-Slurm from
version 0.10.0 onwards, assisting with code generation, refactoring,
test scaffolding, and documentation drafting. It was also used to assist
with drafting and editing this paper. All AI-assisted outputs were
reviewed, edited, and validated by the author, who made all core design
decisions and retains full responsibility for the accuracy, originality,
and correctness of the submitted materials.

\section*{References}\label{references}
\addcontentsline{toc}{section}{References}

\phantomsection\label{refs}
\begin{CSLReferences}{1}{0}
\bibitem[\citeproctext]{ref-di2017nextflow}
Di Tommaso, P., Chatzou, M., Floden, E. W., Barja, P. P., Palumbo, E.,
\& Notredame, C. (2017). Nextflow enables reproducible computational
workflows. \emph{Nature Biotechnology}, \emph{35}(4), 316--319.
\url{https://doi.org/10.1038/nbt.3820}

\bibitem[\citeproctext]{ref-lannelongue2021green}
Lannelongue, L., Grealey, J., \& Inouye, M. (2021). Green algorithms:
Quantifying the carbon footprint of computation. \emph{Advanced
Science}, \emph{8}(12). \url{https://doi.org/10.1002/advs.202100707}

\bibitem[\citeproctext]{ref-molder2021snakemake}
Mölder, F., Jablonski, K. P., Letcher, B., Hall, M. B., Tomkins-Tinch,
C. H., Sochat, V., Forster, J., Lee, S., Twardziok, S. O., Kanitz, A.,
Wilm, A., Holtgrewe, M., Rahmann, S., Nahnsen, S., \& Köster, J. (2021).
Sustainable data analysis with snakemake. \emph{F1000Research},
\emph{10}, 33. \url{https://doi.org/10.12688/f1000research.29032.2}

\bibitem[\citeproctext]{ref-slurm_adoption}
Wang, B., Chen, Z., \& Xiao, N. (2020). A survey of system scheduling
for HPC and big data. \emph{Proceedings of the 2020 4th International
Conference on High Performance Compilation, Computing and
Communications}, 178--183. \url{https://doi.org/10.1145/3407947.3407977}

\bibitem[\citeproctext]{ref-jette2002slurm}
Yoo, A. B., Jette, M. A., \& Grondona, M. (2003). SLURM: Simple linux
utility for resource management. \emph{Job Scheduling Strategies for
Parallel Processing}, 44--60. \url{https://doi.org/10.1007/10968987_3}

\end{CSLReferences}

\end{document}
